\chapter{Implementation and Experimental Protocol}
\label{ch:implementation}

% =============================================================================
% SKELETON. Written once the model in simulation/ exists.
% Target: ~2,000 words. Every section below states what it must contain.
% Chapter 4 is the specification; this chapter is how that specification was
% realised in code, and how the experiments were run. Do not re-specify rules
% here --- cross-reference the submodel instead.
% =============================================================================

\section{Framework and Architecture}
\label{sec:implementation-architecture}

% CONTENT REQUIRED:
% - Mesa version, Python version, repository layout under simulation/.
% - Class structure: HouseholdAgent, TraditionalLender, BNPLPlatform, CreditBureau,
%   and the Model class that owns the scheduler and the data collector.
% - Which ODD entity (Section~\ref{sec:entities}) maps to which class.
% - How the population parquet from Chapter~\ref{ch:data} is loaded and instantiated (P5).
% - A figure: class/interaction diagram. This is the natural place for the first
%   structural figure in the thesis.

\section{The Tick}
\label{sec:implementation-tick}

% CONTENT REQUIRED:
% - The seven-step ordering of Section~\ref{sec:scheduling} as implemented, step by step.
% - Where the one-tick lag on the reference-group adoption share is stored and when it is
%   recomputed. This is the detail that makes the peer channel activation-order independent,
%   so it is worth showing concretely.
% - How random asynchronous activation is reseeded each tick (Submodel~17).

\section{Parameter Register}
\label{sec:parameters}

% CONTENT REQUIRED:
% - One table: parameter, symbol, baseline value, sweep range, source or "assumption".
% - Must agree exactly with Table~\ref{tab:submodels} and with
%   Appendix~\ref{app:specification}. Generate it from the config file rather than
%   typing it, so the three cannot drift apart.
% - Flag the two uncalibrated parameters (q_base, beta) and the one unanchored rule
%   (Submodel~4) in the table itself, not only in the text.

\section{Experimental Protocol}
\label{sec:protocol}

% CONTENT REQUIRED:
% - Burn-in: 12 ticks discarded of 52. State what is checked to confirm burn-in is adequate.
% - Replicates per parameter combination, and the seed scheme. Section~\ref{sec:peer-influence}
%   notes that positive feedback amplifies early random differences, so justify the replicate
%   count with a dispersion argument rather than picking a round number.
% - The full sweep grid: BNPL access rate x platform count (1-6) x beta (including 0)
%   x the four intervention levers x the Submodel~4 amount-rule variants.
% - Total run count and wall-clock cost.
% - The calibration procedure for the income-shock probability: what was searched, over what
%   range, against which statistic of pattern~2, and what the fitted value came out at.
%   This is the step that makes the baseline calibrated rather than validated
%   (Section~\ref{sec:lim-calibration}), so it must be transparent.

\section{Verification}
\label{sec:verification}

% CONTENT REQUIRED:
% - Verification is "did I build the model I specified", as distinct from validation,
%   which is "does the model behave like the world". Keep them separate.
% - Unit tests: balance-sheet conservation (money is neither created nor destroyed),
%   arrears ageing across CCMR bands, Reg 23A gate reproducing both published worked
%   examples, Pay-in-4 schedule landing on tick boundaries, late fee capping at R285.
% - Degenerate-case checks: beta=0 reproduces the independent-agent model exactly;
%   BNPL disabled reproduces the baseline exactly; a single platform reduces stacking
%   to single-platform accumulation.
% - Confirm the peer channel is inert in the baseline (s_g identically zero), since the
%   baseline calibration depends on it.

\section{Reproducibility}
\label{sec:reproducibility}

% CONTENT REQUIRED:
% - Seeds, environment capture, and where raw run output lands (results/raw) versus
%   summarised output (results/summary).
% - How a reader could regenerate every figure in Chapter~\ref{ch:results} from the
%   repository.
